\section{An integer cover for a shared polylogarithmic core}
\label{sc:section}

The preceding fractional argument has an integer consequence for cores
of size \(o(\log\log M)\). A different construction permits a much
larger common core. It uses the source factorization itself: outside
prime powers are covered first, regular core parts are assigned to
disjoint target bins, and the remaining core demands are repaired
globally. No separation between prime sizes is assumed.

\begin{theorem}[A shared polylogarithmic core]
\label{thm:sc-polylog}
Fix \(0<a<2\). For every \(\varepsilon>0\), all sufficiently large
\(M\) have the following property. Suppose
\(A\subseteq\{2,\ldots,M\}\) contains \(n\) distinct integers and
admits a designated prime set \(E\) with
\[
 |E|\le(\log M)^a
\]
such that each source has at most one distinct prime factor outside
\(E\). Arbitrary powers of that outside prime are allowed.
Then, for every positive interval \(I\) of \(M\) consecutive
integers,
\begin{equation}
 \tau(A,I)\le(2+\varepsilon)n.
 \label{sc:main-bound}
\end{equation}
The threshold is uniform in the actual core and outside primes, the
source set, its cardinality, and the interval shift. In particular,
\(\tau(A,I)\le(2+o(1))n\) uniformly over this class as \(n\to\infty\).
\end{theorem}

The core may vary with the instance. It is shared by the entire
source set: a mixed source has the unique form \(d r^e\), where
\(d\) is \(E\)-smooth, \(r\notin E\) is prime, and \(e\ge1\).
The proof below is entirely integral. It uses neither a minimal
counterexample nor fractional rounding.

\subsection{One target per outside-prime occurrence group}

\begin{lemma}[Nested prime-power matching]
\label{sc:outside-matching}
For a distinct source set \(A\subseteq[2,M]\) and a prime \(r\),
the source exponents \(v_r(a)>0\) can be assigned to distinct
targets \(b\in I\), each satisfying \(v_r(b)\ge v_r(a)\).
Consequently, when each source has at most one distinct prime
outside \(E\), there is a target set \(B_{\rm out}\) of size at
most \(n_{\rm out}\) covering the full demand of every outside
prime. Here \(n_{\rm out}\) counts all mixed sources.
\end{lemma}

\begin{proof}
For every \(j\ge1\), distinctness gives
\begin{equation}
 \#\{a\in A:v_r(a)\ge j\}
 \le\left\lfloor\frac M{r^j}\right\rfloor
 \le\#\{b\in I:r^j\mid b\}.
 \label{sc:nested-hall}
\end{equation}
The allowed target sets are nested. Order the source exponents
decreasingly and assign them in that order. At a current exponent
\(e\), its allowed target set has at least as many positions as
there are source exponents at least \(e\). The previous assignments
have used fewer positions than that number, so an allowed unused
target exists. This proves the assertion for one prime.

Perform this matching separately for each outside prime \(r\), and
let \(B_r\) be the selected set. Take
\[
 B_{\rm out}=\bigcup_{r\notin E}B_r.
\]
For each \(r\), the union retains all of \(B_r\), and therefore
covers the entire \(r\)-demand. Different groups demand disjoint
prime coordinates; overlaps between their selected sets cause no
loss. Since each mixed source belongs to exactly one such group,
\(|B_{\rm out}|\le\sum_r|B_r|=n_{\rm out}\).
\end{proof}

\begin{proposition}[A finite certificate valid for every shift]
\label{sc:finite-certificate}
Under the one-distinct-outside-prime source condition, let \(d(a)\)
be the full \(E\)-smooth part of \(a\), and let
\(m_d=\#\{a\in A:d(a)=d\}\).
Choose any collection \(\mathcal D\) of occupied core parts \(d>1\).
Suppose the number of targets with exact core part \(d\) is at least
\(m_d\), for every shift and every \(d\in\mathcal D\). Then
\begin{equation}
 \tau(A,I)\le n_{\rm out}+\sum_{d\in\mathcal D}m_d
       +\sum_{\substack{a\in A\\d(a)>1,\ d(a)\notin\mathcal D}}
                    \Omega(d(a))
 \label{sc:finite-cover}
\end{equation}
for every shift. In particular, if the right side is at most \(2n\),
this certifies the exact \(2n\) bound for the given source set.
\end{proposition}

\begin{proof}
Use Lemma~\ref{sc:outside-matching} for all outside demands.
Choose \(m_d\) distinct targets from each selected exact-core-part
bin; these bins are disjoint. Their product covers all core demand
from those fibres. In the union of the two target sets, only the
remaining core units can be deficient. Repair them one at a time,
using unused targets guaranteed by full-interval supply. At most
the last sum in \eqref{sc:finite-cover} additions are needed.
Core part one requires no carrier.
\end{proof}

This certificate is finite and explicit. For \(T=M/d\),
inclusion--exclusion supplies the uniform bin lower bound
\[
 L_E(T)=
 \sum_{\substack{S\subseteq E\\|S|\ {\rm even}}}
      \left\lfloor\frac T{q_S}\right\rfloor
 -\sum_{\substack{S\subseteq E\\|S|\ {\rm odd}}}
      \left\lceil\frac T{q_S}\right\rceil,
 \qquad q_S=\prod_{p\in S}p,\quad q_\varnothing=1.
\]
Thus \(L_E(M/d)\ge m_d\) is sufficient. The simpler lower bound
\(T\prod_{p\in E}(1-1/p)-2^{|E|}\) is also valid.
The asymptotic proof below obtains better bin estimates when the
core grows; neither finite test is asserted to be necessary.

\subsection{Uniform estimates from the core cardinality}

For \(0<\sigma<1\), let
\[
 \Psi_E(x)=\#\{d\le x:d\text{ is }E\text{-smooth}\},
 \qquad
 {\cal Z}_E(\sigma)=\prod_{p\in E}(1-p^{-\sigma})^{-1}.
\]
One is included among the smooth integers. Rankin's elementary
inequality gives, for every \(x>0\),
\begin{equation}
 \Psi_E(x)\le x^\sigma{\cal Z}_E(\sigma).
 \label{sc:rankin}
\end{equation}
For \(x<1\) the left side is zero, so the same bound remains valid.

If \(E=\{p_1<\cdots<p_k\}\), then \(p_i\ge i+1\), and
\[
 \begin{split}
 \log{\cal Z}_E(\sigma)
 &\le\frac1{1-2^{-\sigma}}\sum_{p\in E}p^{-\sigma}\\
 &\le C_\sigma k^{1-\sigma}.
 \end{split}
\]
The empty Euler product is one. Fix, from now on,
\begin{equation}
 \max\{0,1-1/a\}<\sigma<\frac12.
 \label{sc:sigma-choice}
\end{equation}
Such a choice exists for every fixed \(0<a<2\). It yields
\begin{equation}
 {\cal Z}_E(\sigma)
 \le\exp\!\bigl(C_\sigma(\log M)^{a(1-\sigma)}\bigr)
 =M^{o(1)}
 \label{sc:euler-uniform}
\end{equation}
uniformly over all core prime values, since \(a(1-\sigma)<1\).

\begin{lemma}[Many targets with an exact core part]
\label{sc:bins}
Fix \(a>0\) and \(0<\theta<1\). There is \(c_a>0\) such that,
for all sufficiently large \(M\), every prime set
\(|E|\le(\log M)^a\) and every real interval \(J\) of length
\(M^\theta\le T\le M\) satisfy
\begin{equation}
 \#\{t\in J:(t,E^\Pi)=1\}
 \ge c_a\,\frac T{\log\log M}.
 \label{sc:bin-density}
\end{equation}
The large-\(M\) threshold may depend on \(a,\theta\); the estimate
is uniform in all prime values and shifts.
\end{lemma}

\begin{proof}
We use the inherited uniform interval sieve \cite{Synthesis},
derived from the fundamental lemma \cite{Lagarias,GranvilleSound}.
Precisely, if \(Q\) is a subset of primes at most \(Y\) and an
interval has real length \(T\ge Y^u\), then its count of integers
coprime to \(Q^\Pi\) differs from
\(T\prod_{q\in Q}(1-1/q)\) by at most \(e(u)\) times that
quantity, where \(e(u)\to0\) uniformly in \(Q\) and the shift.
The explicit error is stated in Proposition~\ref{cg:conditional-input}
and in the sieve input of the separated-support section.
We also use
\(\prod_{p\le Y}(1-1/p)\gg1/\log Y\)
\cite{GranvilleSound,TaoMertens}.

Set
\[
 Y=(\log M)^{2a+3},\qquad
 E_{\rm small}=E\cap[2,Y],\qquad E_{\rm large}=E\setminus E_{\rm small}.
\]
Since \(T\ge M^\theta\), the ratio \(\log T/\log Y\) tends to
infinity uniformly. For large \(M\), the small-prime sieve therefore
supplies at least
\[
 \frac12 T\prod_{q\in E_{\rm small}}(1-1/q)
 \gg\frac T{\log Y}
\]
integers coprime to \(E_{\rm small}\).
Requiring coprimality to the remaining core primes discards at most
\[
 \sum_{q\in E_{\rm large}}\left(\frac Tq+1\right)
 \le\frac{T|E|}{Y}+|E|
 =o(T/\log Y).
\]
Indeed
\[
 \frac{|E|\log Y}{Y}\to0,\qquad
 \frac{|E|\log Y}{T}\to0
\]
uniformly at the stated range of \(T\). Subtracting this upper
bound leaves a fixed positive multiple of \(T/\log Y\).
Finally \(\log Y=(2a+3)\log\log M\), proving the lemma.
\end{proof}

If \(d\) is \(E\)-smooth, the target positions with full
\(E\)-smooth part exactly \(d\) are
\[
 \{b\in I:d\mid b,\ (b/d,E^\Pi)=1\}.
\]
They are counted by applying Lemma~\ref{sc:bins} to the scaled
real interval \(I/d\), of length \(M/d\). Different \(d\)'s
give disjoint target bins.

\subsection{The three-stage construction}

\begin{proof}[Proof of Theorem~\ref{thm:sc-polylog}]
Keep the fixed \(\sigma\) from \eqref{sc:sigma-choice}, and put
\begin{equation}
 \theta=\frac{1/2-\sigma}{2(1-\sigma)},\qquad
 Z=M^\theta,\qquad
 \beta=\sigma+\theta(1-\sigma)=\frac14+\frac{\sigma}{2}<\frac12.
 \label{sc:cutoff}
\end{equation}
The parameters depend only on \(a\).
For every mixed source write \(a=d f\), where \(d\) is its full
\(E\)-smooth part and \(f=r^e\) is its outside prime-power part.
Call it regular when \(M/d\ge Z\), and bad otherwise.
All pure \(E\)-smooth sources are also called bad.
Let \(n_{\rm reg}\) be the number of regular sources, and let
\(n_{\rm out}\) count all mixed sources, including bad ones.

\emph{The bad sources have negligible total core demand.}
For a bad mixed source, \(d>M/Z\) and \(df\le M\), so \(f<Z\).
Enlarge the possible outside prime powers to all integers
\(2\le f<Z\). Equations \eqref{sc:rankin}--\eqref{sc:euler-uniform}
give
\begin{equation}
 \begin{split}
 |A_{\rm bad}|
 &\le\Psi_E(M)+\sum_{2\le f<Z}\Psi_E(M/f)\\
 &\le{\cal Z}_E(\sigma)M^\sigma
                    \left(1+\sum_{2\le f<Z}f^{-\sigma}\right)\\
 &\le C_\sigma{\cal Z}_E(\sigma)M^\sigma Z^{1-\sigma}
 =M^{\beta+o(1)}.
 \end{split}
 \label{sc:bad-count}
\end{equation}
Dropping the restriction \(d>M/Z\) only enlarged the count.
Write \(\Omega_E(a)=\sum_{p\in E}v_p(a)\).
The total bad core demand is at most
\begin{equation}
 \sum_{a\in A_{\rm bad}}\Omega_E(a)
 \le |A_{\rm bad}|\log_2M
 =M^{\beta+o(1)}
 =o(\sqrt M).
 \label{sc:bad-demand}
\end{equation}
All errors are uniform in the admitted source data.

\emph{First select targets for every outside-prime demand.}
Lemma~\ref{sc:outside-matching} gives a set \(B_{\rm out}\)
of at most \(n_{\rm out}\) positions. Its product covers the
entire demand of each prime outside \(E\), including those in
bad mixed sources.

\emph{Next select simultaneous carriers for the regular core parts.}
Group regular sources by their exact core part \(d\). Put
\(T_d=M/d\), and let \(m_d\) be the fibre size.
Distinct sources in the same fibre have distinct outside prime
powers \(f\le T_d\). If \({\cal P}(T)\) counts prime powers
\(p^j\le T\), \(j\ge1\), then
\[
 m_d\le{\cal P}(T_d),\qquad
 {\cal P}(T)\le\pi(T)+
       \sum_{2\le j\le\log_2T}\pi(T^{1/j})
 =O(T/\log T).
\]
Here Chebyshev bounds the first term; the remaining terms total
at most \(\sqrt T\log_2T=o(T/\log T)\).
On every occupied regular fibre \(T_d\ge M^\theta\), so
Lemma~\ref{sc:bins} gives an exact-core-part bin of size
\[
 N_d\ge c_a\,\frac{T_d}{\log\log M}.
\]
Consequently
\[
 \frac{m_d}{N_d}
 =O_a\!\left(\frac{\log\log M}{\log T_d}\right)
 =O_a\!\left(\frac{\log\log M}{\log M}\right)\longrightarrow0.
\]
Thus \(N_d\ge m_d\) simultaneously in all fibres for large
\(M\). Choose \(m_d\) distinct targets from each bin. Since the
bins are disjoint, this yields a set \(B_{\rm core}\) of exactly
\(n_{\rm reg}\) positions with
\begin{equation}
 \prod_{a\in A_{\rm reg}}d(a)\ \bigm|\ \prod_{b\in B_{\rm core}}b.
 \label{sc:core-carriers}
\end{equation}
This is one joint covering of all regular core multiplicities.

\emph{Finally repair the remaining core units globally.}
Take \(B_0=B_{\rm out}\cup B_{\rm core}\). Its cardinality is
at most \(n_{\rm out}+n_{\rm reg}\).
For every outside prime it retains all of \(B_{\rm out}\);
for every core prime it retains all of \(B_{\rm core}\).
Therefore the only possible deficits are core units from the
bad sources, whose total is bounded in \eqref{sc:bad-demand}.
If a prime remains deficient, full-interval supply
\(\prod A\mid M!\mid\prod I\) guarantees an unused target
divisible by that prime. Add such a target. Each addition reduces
the total positive deficit by at least one and never increases
another deficit. Hence
\begin{equation}
 \begin{split}
 \tau(A,I)
 &\le n_{\rm out}+n_{\rm reg}
                 +\sum_{a\in A_{\rm bad}}\Omega_E(a)\\
 &\le2n+M^{\beta+o(1)}.
 \end{split}
 \label{sc:constructive-bound}
\end{equation}
Overlaps between the first two selected sets are harmless, since
their demanded prime families are disjoint. The last repair uses
the accumulated cover, rather than an independent cover of the
bad sources.

To finish uniformly over all cardinalities, if \(n\le\sqrt M/4\),
the ambient quadratic theorem \cite[Appendix A]{Wang} gives
\(\tau(A,I)\le2n\). Otherwise \(n>\sqrt M/4\), and
\[
 \frac{M^{\beta+o(1)}}n
 \le4M^{\beta-1/2+o(1)}\longrightarrow0.
\]
This proves \eqref{sc:main-bound} for all sufficiently large
\(M\), uniformly over the claimed class and shifts.
Finally \(n\to\infty\) forces \(M\ge n+1\to\infty\),
giving the stated uniform asymptotic conclusion.
\end{proof}

\subsection{An explicit core bound approaching the squared logarithm}

\begin{corollary}[A quantified near-square core]
\label{cor:sc-near-square}
For all sufficiently large \(M\), suppose
\begin{equation}
 |E|\le10^{-6}
       \left(\frac{\log M}{\log\log M}\right)^2
 \label{sc:near-square-size}
\end{equation}
and each distinct source has at most one distinct prime outside
\(E\), with arbitrary powers allowed. Then
\(\tau(A,I)\le(2+o(1))|A|\), uniformly in the source set,
cardinality, prime values, and shift. More precisely, with
\(L=\log M\) and \(l=\log L\), the direct construction gives
\begin{equation}
 \tau(A,I)\le2|A|+
 O\!\left(\sqrt M\,L\exp\left(-\frac{2L}{5l}\right)\right).
 \label{sc:near-square-error}
\end{equation}
The absolute numerical constant in \eqref{sc:near-square-size}
is a sufficient choice and is not asserted to be optimal.
\end{corollary}

\begin{proof}
We check the varying parameters in the same construction explicitly,
rather than substituting them into a theorem for fixed parameters.
Assume \(l\ge4\), and set
\[
 \delta=\frac1l,\qquad
 \sigma=\frac12-\delta\in[1/4,1/2),\qquad
 \theta=\frac{\delta}{2(1-\sigma)}=\frac1{l+2},\qquad
 \beta=\frac12-\frac1{2l}.
\]
Let \(k=|E|\). The Euler estimate is uniform in this range of
\(\sigma\):
\[
 \begin{split}
 \log{\cal Z}_E(\sigma)
 &\le\frac1{1-2^{-\sigma}}
                       \sum_{i=1}^k(i+1)^{-\sigma}\\
 &\le\frac{2}{1-2^{-1/4}}\,k^{1-\sigma}
 <13k^{1-\sigma}.
 \end{split}
\]
The empty core is immediate. For \(k\ge1\), the bound \(k\le L^2\)
implies \(k^\delta\le e^2\). Therefore
\begin{equation}
 \log{\cal Z}_E(\sigma)
 <13e^2\sqrt k<100\sqrt k\le\frac{L}{10l}.
 \label{sc:near-square-euler}
\end{equation}
Here \(\sqrt k\le10^{-3}L/l\) follows from
\eqref{sc:near-square-size}.

Define the regular mixed sources by \(M/d\ge M^\theta\), as
before. The sum in \eqref{sc:bad-count} is bounded uniformly by
\(3(M^\theta)^{1-\sigma}\), since \(1/(1-\sigma)\le2\).
Combining this with \eqref{sc:near-square-euler} gives
\[
 |A_{\rm bad}|
 \le3\sqrt M\exp\left(-\frac{L}{2l}+\frac{L}{10l}\right)
 =3\sqrt M\exp\left(-\frac{2L}{5l}\right).
\]
Consequently the total bad core demand is at most this bound times
\(\log_2M\), which is \(o(\sqrt M)\).

It remains to check the target bins with the shrinking \(\theta\).
Take \(Y=L^7\). Every occupied regular fibre has
\[
 T\ge M^\theta=\exp\left(\frac L{l+2}\right),\qquad
 \frac{\log T}{\log Y}\ge\frac{L}{7l(l+2)}\longrightarrow\infty.
\]
The uniform small-prime sieve still gives a fixed positive multiple
of \(T/\log Y\) positions coprime to \(E\cap[2,Y]\).
Deleting multiples of the remaining core primes costs at most
\[
 \frac{Tk}{Y}+k\le\frac T{L^5}+L^2=o(T/\log Y).
\]
Thus exact-core-part bins have size \(\gg T/l\).
The number of possible prime-power remainders remains
\(O(T/\log T)\), and its ratio to the bin size is
\[
 O\left(\frac l{\log T}\right)
 =O\left(\frac{l(l+2)}L\right)\longrightarrow0.
\]
All bins can therefore be assigned as before. The outside-prime
matching and the final global repair do not depend on
\(\sigma,\theta\). They yield \eqref{sc:near-square-error}.

If \(|A|\le\sqrt M/4\), use the exact quadratic theorem.
Otherwise divide the error in \eqref{sc:near-square-error} by
\(|A|>\sqrt M/4\). The relative error is
\(O(L\exp(-2L/(5l)))=o(1)\), uniformly, proving the corollary.
\end{proof}

For every fixed \(a<2\), one has
\((\log M)^a\le10^{-6}(\log M/\log\log M)^2\) eventually.
Thus the quantified corollary includes the fixed-exponent theorem,
while its proof checks the additional uniformity needed when the
effective exponent approaches two.

\begin{corollary}[A stronger core order from Chebyshev's estimate]
\label{cor:sc-chebyshev-core}
There is an absolute constant \(c_0>0\) such that the same uniform
integer conclusion holds under the source factorization hypothesis
of Theorem~\ref{thm:sc-polylog}, whenever
\begin{equation}
 |E|\le c_0\,\frac{(\log M)^2}{\log\log M}.
 \label{sc:chebyshev-core-size}
\end{equation}
For sufficiently large \(M\), the construction also satisfies the
error bound \eqref{sc:near-square-error}.
The existence of a positive absolute \(c_0\) is asserted; no optimal
constant is claimed.
\end{corollary}

\begin{proof}
Chebyshev's estimate implies \(p_i\gg i\log(i+1)\) for the
\(i\)-th prime \cite{GranvilleSound}. Uniformly for
\(1/4\le\sigma\le1/2\), this gives
\begin{equation}
 \log{\cal Z}_E(\sigma)
 \le C\,\frac{k^{1-\sigma}}{(\log(k+1))^\sigma},
 \qquad k=|E|\ge1,
 \label{sc:chebyshev-euler}
\end{equation}
with an absolute \(C\). To check uniformity, split the sum of
\(p_i^{-\sigma}\) at \(i=\sqrt k\). The first part is
\(O(k^{(1-\sigma)/2})\). In the second part
\(\log(i+1)\gg\log(k+1)\), giving
\(O(k^{1-\sigma}/(\log(k+1))^\sigma)\).
The first part is absorbed in the second bound uniformly in
this compact range of \(\sigma\); its relative ratio is at most
a constant times \(\sqrt{\log(k+1)}/k^{1/4}\).
The Euler logarithm introduces only the bounded factor
\((1-2^{-\sigma})^{-1}\). Enlarging \(C\) handles small \(k\).

Use \(L=\log M\), \(l=\log L\), and the same
\(\sigma=1/2-1/l\) as in Corollary~\ref{cor:sc-near-square}.
Put \(K=\lfloor c_0L^2/l\rfloor\).
The function \(k^{1-\sigma}/(\log(k+1))^\sigma\) is increasing:
its logarithmic derivative multiplied by \(k\) is
\[
 (1-\sigma)-
 \frac{\sigma k}{(k+1)\log(k+1)}
 \ge1-2\sigma\ge0.
\]
For every fixed \(c_0>0\), as \(M\to\infty\),
\[
 \begin{split}
 \frac{K^{1-\sigma}}{(\log(K+1))^\sigma}
 &=\sqrt{\frac K{\log(K+1)}}\,
   \exp\left(\frac{\log K+\log\log(K+1)}l\right)\\
 &=\left(e^2\sqrt{\frac{c_0}{2}}+o(1)\right)\frac L l.
 \end{split}
\]
Choose \(c_0\le1\) so small that
\(Ce^2\sqrt{c_0/2}<1/20\). Then
\eqref{sc:chebyshev-euler}, for every \(k\le K\), is at most
\(L/(10l)\) once \(M\) is large.
The empty core causes no error.

This is precisely the Euler bound
\eqref{sc:near-square-euler} needed in the preceding proof.
Also \(k\le L^2\) for large \(M\), so its \(Y=L^7\)
sieve calculation, large-prime deletion bound, prime-power fibre
comparison, and bad-demand estimate remain valid without change.
The same construction and the quadratic gate finish the proof.
\end{proof}

\subsection{Scope of the larger shared-core class}

For example, a designated core of size at most \((\log M)^{1.99}\)
is permitted, with any number of possible outside primes and no
prime-size gap. The source restriction still concerns the number
of \emph{distinct} outside primes, so their arbitrary powers remain
fully included.

The two conditions on the Rankin exponent are
\[
 a(1-\sigma)<1,\qquad \sigma<1/2.
\]
Their simultaneous use explains the endpoint \(a<2\) in the
fixed-parameter argument. Corollaries~\ref{cor:sc-near-square}
and \ref{cor:sc-chebyshev-core} check specified approaches to
that endpoint. These results do not
establish a core of size \((\log M)^2\), an arbitrary
instance-dependent exponent tending to two, or an optimal endpoint.

The unrestricted minimal-obstruction core can still have size
\(O(\sqrt{n/\log n})\), much larger than a fixed power of
\(\log M\) in the subquadratic range. Thus
Theorem~\ref{thm:sc-polylog} enlarges a concrete source class
without resolving the unrestricted asymptotic or exact conjecture.
Its integer construction is independent of the controlled-weight
hinge bound: the latter remains useful for restricting dual vectors
inside a larger designated core.
